%===============================================================================
\section{Discussion}\label{sec:discussion}
%===============================================================================

\paragraph{What changes for grading practice}
Single-field grading is not a conservative approximation of a combination-aware audit; it fails
in the unsafe direction. On every data set, most fields that take part in a dangerous combination of
at most three fields are safe on their own, and withholding what single-field grading flags leaves
more than 95\% of those combinations intact. The field-level table proposed here keeps the
familiar form of a grade per field but adds what a grade needs in order to act on it: the
worst-case risk $\Mik$, the witness background that attains it, the escalation over single-field
risk, and the minimal unsafe sets from which a minimum withholding set follows. Because $\Mik$ is
defined relative to the public baseline, the table is recomputed whenever something is released;
the amortized attacker makes this recomputation cheap and the certification step keeps it honest.

\paragraph{Relation to inference-graph approaches}
Propagating risk over pairwise inference relations~\cite{hu2026ignn} captures chains of
single-field inferences and physical relations known to experts, and it does not require an
attacker to be trained. It cannot represent joint inference by construction (Fig.~\ref{fig:mechanism}),
and its scores are not checkable by an attack experiment. Set-level inferability provides the
missing ground truth; an inference graph could be used to propose candidate backgrounds for the
scan, a combination we have not evaluated.

\paragraph{What the estimator is and is not}
The amortized attacker is not required for correctness: every decision reported here is certified
by retraining. Its role is to make the scan cheap enough to be repeated, and its accuracy matters
only through the ranking of backgrounds that certification follows. The ablation shows that most
of this accuracy comes from solving a readout per field set on a nonlinear conditional
representation; masked pre-training chiefly improves the ordering of the largest marginals, and
one pre-training on all columns serves any later choice of targets and candidates, with
certification required for targets unseen in pre-training.

\paragraph{Limitations}
\emph{(i)} Risk is defined relative to a fixed attacker family; a stronger attacker can only raise
$\Vy$ and hence the measured combination risk, but the numerical values are family-dependent.
\emph{(ii)} The adversary is assumed to hold labelled history of the target; with little auxiliary
data, $\Mik$ as a maximum statistic is inflated by selection noise, which again argues for certified
lower bounds rather than point values. \emph{(iii)} We use random train/test splits; drift over time
may change the values of $\Mik$, although not the audit procedure. \emph{(iv)} $K=2$ nearly saturates the marginal risk on nine of ten targets but not threshold
crossings; disposal should use $K=3$, and parity-like targets would defeat any fixed $K$. \emph{(v)} Field roles for NEM, PJM and CAISO follow the Chinese disclosure rules rather
than local market rules; they are used to define a realistic audit, not to assess those markets.
\emph{(vi)} The inference-graph baseline reimplements the propagation idea without the expert inputs
and attention training of~\cite{hu2026ignn}.

%===============================================================================
\section{Conclusion}\label{sec:conclusion}
%===============================================================================

We defined field-level inference risk from set-level inferability, the out-of-sample accuracy of
retrained attackers on a set of fields. The budgeted marginal risk $\Mik$ is the worst-case gain
that releasing a field brings over any background of at most $K$ fields, relative to what is
already public; it is a threshold-free safety margin, and its critical fields are exactly the
members of small minimal unsafe sets. An amortized attacker---masked pre-training with a
closed-form readout per field set---scans all small sets in milliseconds, and retraining a few
scanned backgrounds per field certifies the result. On four power data sets with rule-based field
roles and exhaustive retraining as ground truth, single-field grading recovered 13--17\% of the
critical fields and left more than 95\% of dangerous combinations intact, whereas the scan--certify
pipeline recovered 75--88\% of critical fields with no false positives and, selecting fields from
estimates alone, broke all but 1--6\% of dangerous combinations with close to the minimum number of
withheld fields.

\section*{Data and code availability}
RTS-GMLC~\cite{barrows2020rts} and the AEMO NEMWEB archive~\cite{aemo2024nemweb} are public; the PJM and CAISO series
were obtained from the operators' public data portals. Field-role tables, the code for dispatch simulation, ground-truth
retraining, the amortized attacker and all analyses will be made available upon publication.
